%%%%%%%%%%%%%%%%%%%%%%%%%%%%%%%%%%%%%%%%%%%%%%%%%%%%%%%%%%%%%%%%%%%%%%%%%%%
% Englische Fassung des Uebungsblatts zu Lehreinheit 6 (Teil A und B).
% Vorlage: handouts/LE06.tex. Nur in handouts.tex mit \def\wuzsprache{en}.
% Datensatz (Task 3) identisch mit der deutschen Fassung; Spalten
% umbenannt: St -> W (wonder), VW -> PK (prior knowledge).
% Da es kein englisches Handout gibt, steht die Faustformel fuer den
% Bayes-Faktor direkt in Task 3.
%%%%%%%%%%%%%%%%%%%%%%%%%%%%%%%%%%%%%%%%%%%%%%%%%%%%%%%%%%%%%%%%%%%%%%%%%%%
\begin{uebungsblatt}{Cups, Dice and Data}

\aufgabe{Monty Hall in pairs}[in class, pairs, 10 min.]
The host secretly rolls a die to place the ball (1--2: A, 3--4: B,
5--6: C), the player points at a cup, and the host lifts an \emph{empty}
cup that was not chosen. Odd rounds: stay; even rounds: switch. Tick
whether the round was won.

\begin{center}
\small
\renewcommand{\arraystretch}{1.25}
\begin{tabular}{|c|c|c|c||c|c|c|c|}
\hline
\textbf{Rd.} & \textbf{Strategy} & \textbf{Ball} & \textbf{won} &
\textbf{Rd.} & \textbf{Strategy} & \textbf{Ball} & \textbf{won} \\ \hline
1 & stay & & & 13 & stay & & \\ \hline
2 & switch & & & 14 & switch & & \\ \hline
3 & stay & & & 15 & stay & & \\ \hline
4 & switch & & & 16 & switch & & \\ \hline
5 & stay & & & 17 & stay & & \\ \hline
6 & switch & & & 18 & switch & & \\ \hline
7 & stay & & & 19 & stay & & \\ \hline
8 & switch & & & 20 & switch & & \\ \hline
9 & stay & & & 21 & stay & & \\ \hline
10 & switch & & & 22 & switch & & \\ \hline
11 & stay & & & 23 & stay & & \\ \hline
12 & switch & & & 24 & switch & & \\ \hline
\end{tabular}
\end{center}
Total: \textbf{stay} won \rule{12mm}{0.3pt} of \rule{12mm}{0.3pt}
\qquad \textbf{switch} won \rule{12mm}{0.3pt} of \rule{12mm}{0.3pt}
\qquad whole class (board): \rule{25mm}{0.3pt}

\aufgabe{Dice detectives}[in class, pairs, 10 min.]
The roller opens the envelope and reads the card \emph{privately}:
\enquote{fair} (only a 6 counts as a hit) or \enquote{loaded} (5 and 6
count as hits). The roller throws the die 30 times, hidden under a cup,
and only says \enquote{hit} or \enquote{no hit}. The detective keeps the
tally. Known in advance: one envelope in five is loaded.

\begin{center}
\renewcommand{\arraystretch}{1.6}
\begin{tabular}{|p{0.3\textwidth}|p{0.55\textwidth}|}
\hline
Hit (tally) & \\ \hline
No hit (tally) & \\ \hline
\end{tabular}
\end{center}

\begin{center}
\footnotesize
\setlength{\tabcolsep}{3.5pt}
\begin{tabular}{@{}l*{12}{c}@{}}
\toprule
Hits $k$ & 3 & 4 & 5 & 6 & 7 & 8 & 9 & 10 & 11 & 12 & 13 & 14 \\
\midrule
$p=P(X\geq k\mid\text{fair})$ & 0.90 & 0.76 & 0.58 & 0.38 & 0.22 & 0.11 & 0.051 & 0.020 & 0.007 & 0.002 & 0.0005 & 0.0001 \\
$\mathit{BF}_{10}=2^k\cdot0.8^{30-k}$ & 0.02 & 0.05 & 0.12 & 0.30 & 0.76 & 1.9 & 4.7 & 11.8 & 30 & 74 & 184 & 461 \\
$P(\text{loaded})$ at $1:4$ & 0.5\,\% & 1\,\% & 3\,\% & 7\,\% & 16\,\% & 32\,\% & 54\,\% & 75\,\% & 88\,\% & 95\,\% & 98\,\% & 99\,\% \\
\bottomrule
\end{tabular}
\end{center}

\begin{enumerate}
  \item \textbf{Frequentist:} $H_0$: fair. p-value for your $k$:
    \rule{20mm}{0.3pt}. At $\alpha=0.05$ the die is
    $\square$ significantly loaded \quad $\square$ not significant.
  \item \textbf{Bayesian:} $\mathit{BF}_{10}$ = \rule{18mm}{0.3pt};
    prior odds $1:4$, posterior odds \rule{18mm}{0.3pt};\\
    $P(\text{loaded})$ = \rule{18mm}{0.3pt}. Do the calculation once
    yourself, without the last row of the table.
  \item Open the envelope. Who was right? Enter your pair in the
    two-by-two table on the board.
  \item What would your Bayesian verdict be if only one die in 100 were
    loaded? And your test verdict?
\end{enumerate}
\schreiblinien{2}

\aufgabe{Analysing a data set}[homework, small groups]
\textbf{Fictitious data set} (constructed for this exercise, not a real
survey): 50 people held an ideomotor pendulum (cf.\ Unit~1). The largest
swing was measured in millimetres (\textbf{mm}); afterwards the people
rated how amazed they were (\textbf{W} for wonder, 1 = not at all to
5 = very much) and stated whether they had known the pendulum effect
before (\textbf{PK} for prior knowledge, 1 = yes, 0 = no).

\begin{center}
\footnotesize
\setlength{\tabcolsep}{3.2pt}
\begin{tabular}{@{}rrrr@{\hspace{9pt}}rrrr@{\hspace{9pt}}rrrr@{\hspace{9pt}}rrrr@{\hspace{9pt}}rrrr@{}}
\toprule
\textbf{No} & \textbf{mm} & \textbf{W} & \textbf{PK} &
\textbf{No} & \textbf{mm} & \textbf{W} & \textbf{PK} &
\textbf{No} & \textbf{mm} & \textbf{W} & \textbf{PK} &
\textbf{No} & \textbf{mm} & \textbf{W} & \textbf{PK} &
\textbf{No} & \textbf{mm} & \textbf{W} & \textbf{PK} \\
\midrule
 1 & 11 & 3 & 0 & 11 &  9 & 3 & 0 & 21 & 10 & 2 & 1 & 31 & 43 & 3 & 1 & 41 & 47 & 4 & 0 \\
 2 &  5 & 2 & 0 & 12 & 24 & 2 & 0 & 22 & 26 & 4 & 1 & 32 & 57 & 4 & 0 & 42 & 59 & 3 & 0 \\
 3 & 15 & 3 & 1 & 13 & 24 & 4 & 1 & 23 & 39 & 4 & 0 & 33 & 21 & 3 & 0 & 43 & 26 & 3 & 0 \\
 4 & 41 & 4 & 1 & 14 &  8 & 2 & 0 & 24 & 15 & 4 & 0 & 34 & 96 & 5 & 0 & 44 & 31 & 4 & 1 \\
 5 & 37 & 3 & 0 & 15 & 16 & 2 & 1 & 25 & 13 & 3 & 0 & 35 & 49 & 4 & 0 & 45 & 23 & 2 & 0 \\
 6 & 39 & 5 & 0 & 16 & 25 & 4 & 0 & 26 & 13 & 3 & 0 & 36 & 30 & 4 & 0 & 46 & 20 & 3 & 1 \\
 7 & 18 & 3 & 0 & 17 & 19 & 5 & 0 & 27 & 66 & 5 & 0 & 37 & 23 & 4 & 0 & 47 & 12 & 1 & 1 \\
 8 &  0 & 1 & 0 & 18 & 53 & 4 & 1 & 28 & 35 & 3 & 1 & 38 & 53 & 4 & 0 & 48 & 36 & 4 & 1 \\
 9 & 16 & 2 & 1 & 19 &  5 & 2 & 1 & 29 & 38 & 5 & 0 & 39 & 18 & 2 & 0 & 49 & 18 & 3 & 0 \\
10 & 31 & 4 & 1 & 20 & 45 & 5 & 0 & 30 & 26 & 3 & 1 & 40 & 73 & 5 & 0 & 50 & 40 & 4 & 0 \\
\bottomrule
\end{tabular}
\end{center}

\begin{enumerate}
  \item State the level of measurement of each variable.
  \item For \textbf{mm}, compute the minimum, maximum, range, median,
    $Q_1$, $Q_3$, $\mathit{IQR}$ and mean. State how you determined the
    quartiles.
  \item Draw a histogram (bin width 10\,mm) and a box plot. Are there
    outliers according to the $1.5\cdot\mathit{IQR}$ rule (values more
    than $1.5\cdot\mathit{IQR}$ below $Q_1$ or above $Q_3$)? What would you
    do with them?
  \item Compute the mean and median of \textbf{mm} without the outlier.
    Which statistic reacts more strongly, and why?
  \item Summarise \textbf{W} in a sensible way. Is a mean admissible?
  \item Compare \textbf{mm} between people with and without prior
    knowledge, using the mean \emph{and} the median. What do you notice?
  \item For the same comparison, compute a $t$-test with the pooled
    standard deviation and the Bayes factor from the rule of thumb
    \parencite{wagenmakers2007practical}, where $N$ is the total number of
    people:
    \[t=\frac{\bar{x}_1-\bar{x}_2}{s_p\sqrt{1/n_1+1/n_2}},\quad
      \mathit{df}=n_1+n_2-2,\qquad
      \mathit{BF}_{01}\approx\sqrt{N}\,\Bigl(1+\frac{t^2}{\mathit{df}}\Bigr)^{-N/2}\] What do the two results say, and what do they not say?
\end{enumerate}
\schreiblinien{6}
\end{uebungsblatt}

\begin{uebungsblattb}{4}{Base Rates, Credibility, Your Project}

\aufgabe{Base rate: the taxi and the witness}[homework]
In a city, 85\,\% of taxis are green and 15\,\% are blue. After a
hit-and-run accident, a witness says the taxi was blue. Tests show that
under the visibility conditions of that night he identifies the colour
correctly in 80\,\% of cases \parencite{tversky1982evidential}.
\begin{enumerate}
  \item First, a gut estimate: how likely was the taxi blue?
  \item Calculate with natural frequencies (100 accidents) and draw the
    tree as on the slide about the medical test.
  \item Do the same calculation in odds form: prior odds, the witness's
    Bayes factor, posterior odds.
  \item A second, independent witness with the same accuracy also says
    \enquote{blue}. How likely is it now?
  \item In another city, half of the taxis are blue. With one witness,
    what changes and what does not?
\end{enumerate}
\schreiblinien{5}

\aufgabe{Credibility check with a search engine}[in class, small groups, 12 min.]
Claim (card): \rule{0.7\textwidth}{0.3pt}

Base rate \emph{before} searching: of 100 claims of this kind, about
\rule{12mm}{0.3pt} are true, because\\ \rule{\textwidth}{0.3pt}

\begin{center}
\small
\renewcommand{\arraystretch}{1.7}
\begin{tabular}{|p{0.38\textwidth}|p{0.3\textwidth}|c|c|}
\hline
\textbf{Finding} (source, what does it say?) & \textbf{expected if true? if false?} & \textbf{Factor} & \textbf{Odds} \\ \hline
Start: base rate & & -- & \hspace{12mm} \\ \hline
 & & & \\ \hline
 & & & \\ \hline
 & & & \\ \hline
 & & & \\ \hline
\end{tabular}
\end{center}
Rough factors: $\times10$, $\times3$, $\times1$, $\times\tfrac13$,
$\times\tfrac1{10}$. Only multiply independent findings; ten sites
copying the same press release are \emph{one} finding.
Result: $P(\text{true})\approx$ \rule{15mm}{0.3pt}. Which finding moved
you the most?
\schreiblinien{2}

\aufgabe{Reading a news report}[homework]
\textbf{Fictitious report} (written for this exercise):

\begin{quote}\small
\emph{Magic tricks make you smarter -- now proven.} People who watch a
magic trick before studying remember more. This is shown by an online
survey of 112 students. The researchers recorded 23 aspects of learning
behaviour. On the item \enquote{I was able to concentrate well}, the group
that saw a magic trick came out significantly ahead ($p = 0.04$).
\enquote{This scientifically proves the value of magic in the lecture
hall}, says the head of the study.
\end{quote}
Answer: What was measured, and in whom? What is claimed? Is an effect
size reported? How many comparisons were made, and what does that mean
for the one significant result? What base rate would you assign to the
claim? If you like, apply the same questions to a real article that
presents a study with $p<0.05$ as \enquote{proof}.
\schreiblinien{4}

\aufgabe{Diamond and map for your idea}[homework, A3, bring to Unit~7]
Take your hypothesis from Unit~5 as the \emph{idea} and fill in the
diamond. At each position, note which observation it \emph{allows} and
which it \emph{rules out}.

\begin{center}
\begin{tikzpicture}[font=\sffamily\footnotesize,
  dia/.style={draw=rwuviolet,thick,rounded corners=2pt,
    text width=52mm,minimum height=24mm,align=left,anchor=center}]
  \node[dia] (a) at (-4.4,0) {\textbf{Idea}\\[18mm]};
  \node[dia] (b) at (4.4,0) {\textbf{Opposite}\\[18mm]};
  \node[dia] (c) at (0,3.0) {\textbf{Both}\\[18mm]};
  \node[dia] (d) at (0,-3.0) {\textbf{Neither}\\[18mm]};
  \draw[rwuviolet,thick] (a) -- (c) -- (b) -- (d) -- (a);
\end{tikzpicture}
\end{center}
\begin{enumerate}
  \item Which observation separates the idea from its opposite? Which
    separates the idea from the \enquote{neither} position? Roughly
    estimate how likely the observation is under each position and
    compute the Bayes factor.
  \item Write down the new, sharpened statement.
  \item On an A3 sheet, draw the map of your study (legend on the slide):
    starting village, paths, fog, dragons, bridges, dead ends, blank
    spots, destination. Every element carries a source or a number.
\end{enumerate}
\schreiblinien{3}

\aufgabe{Analysis plan for your project}[homework, basis for the design sketch in Unit~7]
\begin{center}
\renewcommand{\arraystretch}{1.5}
\begin{tabular}{|p{0.38\textwidth}|p{0.52\textwidth}|}
\hline
H\textsubscript{1} (from Unit~5) & \\[6mm] \hline
Outcome variable and level of measurement & \\ \hline
Comparison: two groups, repeated measures or proportions? & \\ \hline
Smallest effect that would interest you & \\ \hline
Realistic $n$ (who, how many?) & \\ \hline
Power according to the table on the slide & \\ \hline
Analysis: test, estimate with interval, or Bayes factor with stopping rule? & \\ \hline
Base rate of your hypothesis: how surprising would it be if it were true? & \\ \hline
Conclusion: does the question fit the $n$, or will you change the question or the design? & \\[6mm] \hline
\end{tabular}
\end{center}

\aufgabe{Diary of wonder}[homework]
When did your sense of probability last deceive you? Write down what you
expected, what happened, and the matching concept from this unit.

\begin{lehrhinweis}[Expected answers]
\textbf{Task~1.} Per pair, roughly 8 of 12 wins are expected when
switching and 4 of 12 when staying. Individual pairs vary a lot (the
standard deviation of the number of wins in 12 switch rounds is about
1.6); some pairs even find the opposite. Only the class pool shows the
pattern reliably. This is the starting point for talking about sample
size.

\textbf{Task~2.} Table values computed exactly with $B(30,\tfrac16)$ and
$B(30,\tfrac13)$. A die is significant ($\alpha=0.05$, one-sided) from
10 hits on. A fair die gets there with 2.0\,\%, a loaded one with
57\,\%. The Bayesian rule \enquote{loaded if $P(\text{loaded})>50\,\%$}
applies at base rate $1:4$ from 9 hits on (fair: 5.1\,\%, loaded:
71\,\%). Check for $k=10$: $2^{10}=1024$, $0.8^{20}=0.0115$,
$\mathit{BF}_{10}=11.8$; odds $11.8:4$, i.e.\ $11.8/15.8=75\,\%$.
(4)~At base rate $1:99$ and $k=10$: $11.8:99$, i.e.\ 11\,\%. Despite
being \enquote{significant}, the die is then probably fair; the test
verdict does not change because the test does not know the base rate.
With 15 pairs the two-by-two table varies a lot; expected are 0--1 false
alarms and 1--2 of 3 loaded dice detected.

\textbf{Task~3} (values checked with Python).
(1)~No: nominal (identifier); mm: ratio scale; W: ordinal; PK: nominal
(dichotomous).
(2)~$n=50$; minimum 0, maximum 96, range 96\,mm; median 25.5\,mm; mean
29.94\,mm; standard deviation $s=19.32$\,mm. Quartiles as the median of
the lower and upper half: $Q_1=16$, $Q_3=40$, $\mathit{IQR}=24$. With
linear interpolation (e.g.\ Excel \texttt{QUARTILE.INC}) one gets
$Q_3=39.75$ and $\mathit{IQR}=23.75$; both are correct if the method is
stated. Mode: 18 and 26 (three times each).
(3)~Bin frequencies 0--9: 5; 10--19: 13; 20--29: 10; 30--39: 9; 40--49: 6;
50--59: 4; 60--69: 1; 70--79: 1; 80--89: 0; 90--99: 1. The distribution
is right-skewed. Fences: $16-36=-20$ and $40+36=76$. The only outlier is
no.~34 with 96\,mm; 73\,mm (no.~40) lies just inside. With interpolation
the fences are $-19.6$ and $75.4$; same result. Good answers do not
silently delete the value but check it (measurement error? someone swung
on purpose?) and report both analyses.
(4)~Without no.~34: mean 28.59\,mm, median 25.0\,mm, $s=16.98$\,mm. The
mean drops by 1.35\,mm, the median by 0.5\,mm, the standard deviation by
2.3\,mm.
(5)~Frequencies of W: 1: 2; 2: 9; 3: 15; 4: 17; 5: 7. Median 3, mode 4.
Strictly speaking a mean (3.36) is not admissible because the distances
between the levels are unknown; anyone who reports it should justify it.
Optional: Spearman's rank correlation between mm and W is $\rho=0.71$.
(6)~With prior knowledge ($n=17$): mean 25.88\,mm, median 26\,mm,
$s=13.06$. Without prior knowledge ($n=33$): mean 32.03\,mm, median
25\,mm, $s=21.75$. The mean suggests a difference of just over 6\,mm,
the median almost none; the group without prior knowledge varies more
and contains the two largest values (73 and 96\,mm). Expected answer:
these data show no robust group difference, and the choice of statistic
determines the story one could tell.
(7)~Pooled standard deviation $s_p=19.3$\,mm,
$\mathit{SE}=s_p\sqrt{1/33+1/17}=5.76$\,mm,
$t=(32.03-25.88)/5.76=1.07$, $\mathit{df}=48$, $p=0.29$; $d=0.32$.
Welch test: $t=1.25$, $p=0.22$. Rule of thumb:
$\mathit{BF}_{01}\approx\sqrt{50}\,(1+1.14/48)^{-25}\approx7.07\cdot0.56\approx3.9$;
JASP (JZS, $r=\sqrt2/2$) gives $\mathit{BF}_{01}=2.1$. Expected
interpretation: no significant difference; the data lean slightly
towards \enquote{no difference} but do not establish it. The outlier
(no.~34) drives the difference in means; the $t$-test compares means and
inherits this sensitivity.

\textbf{Task~4.} (1)~Most people say 80\,\%. (2)~Of 100 accidents, 15
involve blue taxis, of which the witness identifies 12 as blue; of the 85
green ones he calls 17 blue.
$P(\text{blue}\mid\text{\enquote{blue}})=12/29\approx41\,\%$. (3)~Prior
odds $15:85$, $\mathit{BF}=0.8/0.2=4$, posterior $60:85$, i.e.\ 41\,\%.
(4)~$60:85\times4=240:85$, i.e.\ 74\,\%. (5)~Prior $1:1$, posterior
$4:1$, i.e.\ 80\,\%: the gut feeling is only right when the base rates
are equal. The witness's Bayes factor stays at 4; it describes the
witness, not the city.

\textbf{Task~5.} No \enquote{correct} number is expected. Good groups
justify their base rate before searching, treat copied reports as a
single finding, and give the largest factors to findings that would only
be expected if the claim were true (or false): reviews, fact checks with
sources, the original study itself. Adjusting the base rate after the
search would count the same evidence twice.

\textbf{Task~6.} Measured: self-report, not learning performance; online
survey, unclear whether the groups were randomised. 23 aspects: at
$\alpha=0.05$, pure chance is expected to produce just over one
\enquote{significant} result ($23\cdot0.05\approx1.2$); the probability of
at least one is $1-0.95^{23}\approx69\,\%$. No effect size. The base rate
for \enquote{a magic trick improves learning} is open; for \enquote{it
improves exactly the concentration self-report and nothing else} it is
low. Nothing is \enquote{proven}; at best there is a hypothesis for a
preregistered follow-up study (link to Unit~10).

\textbf{Task~7.} A good opposite is a hypothesis of its own with its own
reasoning, not a mere negation. The \enquote{neither} position should
contain at least one confounder that can be controlled in Unit~7. Key
point (1): an observation that all positions expect equally well is
worthless as evidence; good answers notice this themselves. For the map,
the drawing is not graded, but whether every element carries a source or
a number and whether open dragons are honestly left standing.

\textbf{Task~8.} Good plans name a concrete outcome variable, a justified
smallest effect and an $n$ that can actually be recruited. The most
common insight: with 15--20 people per group only a large effect can be
found; sensible responses are repeated measures, a metric rather than a
yes/no outcome, an estimation question instead of a test, or a
sequential Bayes factor with a maximum sample size.
\end{lehrhinweis}
\end{uebungsblattb}
